%=========================================================================
%  FEX1001 -- Física Experimental I -- UDESC/CCT
%  Modelo de relatório da Experiência 1.
%
%  A Experiência 1 não constrói gráfico linear nem faz linearização: o seu
%  objeto é a medida em si -- algarismos significativos, conversão de
%  unidades, erro de escala, dispersão entre equipes e propagação de erro.
%  Por isso tem estrutura própria, distinta do modelo das Experiências 2 a 5.
%
%  Preenchido automaticamente pela página do laboratório.
%  Compila com pdflatex em qualquer distribuição TeX, inclusive no Overleaf.
%=========================================================================
\documentclass[a4paper,11pt]{article}

\usepackage[utf8]{inputenc}
\usepackage[T1]{fontenc}
\usepackage[brazil]{babel}
\usepackage[margin=2.3cm]{geometry}
\usepackage{amsmath}
\usepackage{amssymb}
\usepackage{graphicx}
\usepackage{booktabs}
\usepackage{siunitx}
\usepackage{float}
\usepackage{caption}
\usepackage[hidelinks]{hyperref}

% group-digits=integer é essencial: sem isso o siunitx agrupa também as
% casas decimais e um valor como 9.79061 sai impresso como "9,790.61".
\sisetup{output-decimal-marker={,}, group-separator={\,}, group-digits=integer}
% Separador de milhar: espaço fino, como prescreve o SI. Ponto e vírgula
% são marcadores decimais conforme a região e não devem agrupar dígitos.
% group-digits=integer impede que o agrupamento alcance as casas decimais,
% que faria 9.79061 sair impresso como "9,790.61".
\captionsetup{font=small,labelfont=bf}

% As tabelas desta experiência são largas (quatro escalas x duas grandezas).
\newcommand{\tabelalarga}{\footnotesize\setlength{\tabcolsep}{4pt}}

%-------------------------------------------------------------------------
% DADOS DO RELATÓRIO -- preenchidos pela página do laboratório
%-------------------------------------------------------------------------
\newcommand{\experimento}{Experiência 1}
\newcommand{\subtitulo}{Física Experimental I}
\newcommand{\equipe}{Bancada}
\newcommand{\integrantes}{Nomes da equipe}
\newcommand{\datarelatorio}{\today}

\newcommand{\disciplina}{FEX1001 -- Física Experimental I}
\newcommand{\professor}{Prof. Dr. Rafael C. R. de Lima}
\newcommand{\instituicao}{Universidade do Estado de Santa Catarina --
Centro de Ciências Tecnológicas, Joinville/SC}

%-------------------------------------------------------------------------
\begin{document}

\begin{center}
  {\large\scshape \instituicao}\\[2pt]
  {\small \disciplina \quad|\quad \professor}\\[14pt]
  \rule{\linewidth}{0.5pt}\\[10pt]
  {\LARGE\bfseries \experimento}\\[4pt]
  {\large \subtitulo}\\[10pt]
  \rule{\linewidth}{0.5pt}\\[10pt]
  {\normalsize Equipe (bancada): \equipe}\\[2pt]
  {\normalsize \integrantes}\\[2pt]
  {\small \datarelatorio}
\end{center}

\vspace{6pt}

%=========================================================================
\section{Objetivo}
\input{secoes/objetivo}

%=========================================================================
\section{Medidas diretas}
\input{secoes/medidas_diretas}

\input{tabelas/tabela1_medidas}

\input{tabelas/tabela1_conversoes}

%=========================================================================
\section{Medidas indiretas}
\input{secoes/medidas_indiretas}

\input{tabelas/tabela2_perimetro}

\input{tabelas/tabela2_area}

%=========================================================================
\section{Teoria de erros e incertezas}
\input{secoes/teoria_erros}

\input{tabelas/tabela3_estatistica}

\input{tabelas/tabela3_areas}

%=========================================================================
\section{Análise dos resultados}
\input{secoes/analise}

%=========================================================================
\section{Conclusões}
\input{secoes/conclusoes}

%=========================================================================
\begin{thebibliography}{9}
\bibitem{halliday} HALLIDAY, D.; RESNICK, R.; WALKER, J.
  \emph{Fundamentos de Física} -- Volume 1: Mecânica.
  8.~ed. Rio de Janeiro: LTC.
\bibitem{apostila} Apostila 1 -- Medidas e algarismos significativos. FEX1001, UDESC/CCT.
\end{thebibliography}

\end{document}
